% !PRINTED 290
\weberpage{290}
%<<CONT
\begin{center}
\S.\,95.

\gesperrt{Beziehung der Charakteristik zu den
Schnittpunkten.}
\end{center}

Man hat nun zu beachten, dass die Bezeichnung der Schnittpunkte
zweier Curven als Ein- oder Austrittsstellen in \S.\,93
wesentlich verschieden ist von der in \S.\,94. Dort kam nur
die Beziehung der Punkte zu den beiden sich schneidenden
Curven in Betracht, während in \S.\,94 noch die Beziehung zu
einer dritten Curve in Frage kam; dies ist in den gewählten
Bezeichnungen $E(\varphi, \psi)$, $A(\varphi, \psi)$ und

\weberfigureimage{weber-b1-fig10.png}{0.48}

$E(f; \varphi, \psi)$, $A(f; \varphi, \psi)$ vollständig ausgedrückt. Nun müssen
wir aber genauer untersuchen, wie sich die beiden Bezeichnungsweisen
zu einander verhalten.

Wir wählen der Deutlichkeit halber eine etwas einfachere
Figur als oben, die zugleich alle möglichen Verhältnisse veranschaulicht
(Fig.~10).

In einem Punkte $E(\varphi, \psi)$ tritt die $\varphi$-Curve in die Fläche
der Curve $\psi$ ein, d.\,h. es geht $\psi$ von positiven zu negativen
Werthen. Ist dann zugleich $f$ positiv, so ist dieser Punkt ein
$E(\varphi; \psi, f)$, wie der Punkt 1 in unserer Figur, ist aber $f$ negativ,
so ist es ein $A(\varphi; \psi, f)$, wie der Punkt 3. So haben wir:

\[
\begin{array}{ll}
\text{Ein Punkt } E(\varphi, \psi) \text{ ist ein} & \\
\quad E(\varphi; \psi, f), \text{ wenn } f > 0, & \text{(Punkt 1 der Fig. 10)} \\
\quad A(\varphi; \psi, f), \text{ wenn } f < 0, & \text{(Punkt 3 der Fig. 10).} \\
\text{Ein Punkt } A(\varphi, \psi) \text{ ist ein} & \\
\quad E(\varphi; \psi, f), \text{ wenn } f < 0, & \text{(Punkt 2 der Fig. 10)} \\
\quad A(\varphi; \psi, f), \text{ wenn } f > 0, & \text{(Punkt 4 der Fig. 10).}
\end{array}
\]

Nun ist in den Punkten $E(\varphi, \psi)$ die Functionaldeterminante
$[\varphi, \psi]$ negativ, in $A(\varphi, \psi)$ positiv (\S.\,93), und wir können
also das Bewiesene so zusammenfassen:

Es ist $E(\varphi; \psi, f)$ die Anzahl der Schnittpunkte von $\varphi$, $\psi$, in
denen $[\varphi, \psi]\, f < 0$,

$A(\varphi; \psi, f)$ die Anzahl der Schnittpunkte von $\varphi$, $\psi$, in
denen $[\varphi, \psi]\, f > 0$;
%CONT>>
